% Source: https://artofproblemsolving.com/community/c3h2013670_geometry
% Stable identity: aops-topic-2013670-q01
The perimeter is $y=2+3+4=9$, so the semiperimeter is $s=\dfrac{9}{2}$.

\begin{asy}
usepackage("amsmath");
size(6cm);
defaultpen(fontsize(10pt));
pair B=(0,0), C=(4,0), A=(11/8,3*sqrt(15)/8);
draw(A--B--C--cycle,linewidth(0.9));
draw(A--(A.x,0),dashed);
dot(A);dot(B);dot(C);
label("$A$",A,N);label("$B$",B,SW);label("$C$",C,SE);
label("$2$",(A+B)/2,NW);label("$3$",(A+C)/2,NE);label("$4$",(B+C)/2,S);
label("\mbox{}",(min(currentpicture,user=true)+max(currentpicture,user=true))/2,black); // keeps the text after the figure black when LaTeX typesets the labels inline
\end{asy}

By Heron's formula,
\[x=\sqrt{s(s-2)(s-3)(s-4)}=\sqrt{\dfrac{9}{2}\cdot\dfrac{5}{2}\cdot\dfrac{3}{2}\cdot\dfrac{1}{2}}=\sqrt{\dfrac{135}{16}}=\dfrac{3\sqrt{15}}{4}.\]
(Check with coordinates: with $B=(0,0)$, $C=(4,0)$, $AB=2$, $AC=3$, the height of $A$ above $BC$ is $\dfrac{3\sqrt{15}}{8}$, and $\dfrac{1}{2}\cdot 4\cdot\dfrac{3\sqrt{15}}{8}=\dfrac{3\sqrt{15}}{4}$.)

Therefore
\[xy=\dfrac{3\sqrt{15}}{4}\cdot 9=\boxed{\dfrac{27\sqrt{15}}{4}}.\]
